\chapter{Des mesures aux indicateurs : construction et validation}
\label{ch:indicateurs}

\questionchapitre{Une mesure brute ne dit rien de la situation qu'elle décrit. Comment produire
des indicateurs normalisés, comparables entre stations et entre années, et comment établir
qu'ils sont justes ?}

\section{Pourquoi standardiser}

Un niveau piézométrique de \SI{12,4}{\metre} n'est ni haut ni bas dans l'absolu : il l'est
relativement à ce que l'ouvrage observe habituellement à cette période de l'année. La même
remarque vaut pour un cumul de pluie ou une température. L'usage opérationnel — décider s'il
faut restreindre les prélèvements, comparer deux bassins, situer un épisode dans l'histoire
climatique — suppose donc de rapporter chaque mesure à sa \textbf{distribution de référence}.

C'est l'objet des indices standardisés. Leur principe est constant : ajuster une loi de
probabilité sur la série de référence, puis transformer la valeur observée en un score exprimé
en écarts-types d'une loi normale centrée réduite. Un score de $-2$ signifie alors la même chose
partout : une situation dont la rareté correspond au seuil de sécheresse sévère.

Deux familles ont été produites : les indices \emph{par station}, calculés sur les niveaux de
nappe et les débits, et les indices \emph{sur grille}, calculés sur les variables climatiques
pour chaque maille du territoire.

\section{Les indices par station}

L'indicateur piézométrique standardisé et son équivalent pour les débits situent le mois courant
par rapport à sa normale saisonnière. La méthode est centralisée dans un unique module, et
l'application ne la recalcule pas : elle lit les tables produites par l'entrepôt.

\begin{center}
\begin{tabularx}{\textwidth}{@{}lX@{}}
\toprule
\textbf{Étape} & \textbf{Détail} \\
\midrule
Référence & Fenêtre fixe 1991--2020, normale climatique retenue par l'Organisation
météorologique mondiale et par le BRGM. Ce n'est pas une fenêtre glissante. \\
Grille & Par mois calendaire, percentiles empiriques de 1 à 99 sur la fenêtre. Un mois comptant
moins de dix observations est interpolé depuis son voisin circulaire. \\
Score & Rang percentile dans la grille, borné à $[0{,}001\,;\,0{,}999]$, puis projeté sur la loi
normale par sa fonction quantile. \\
Classes & Sept classes délimitées par les seuils $z \in \{-1{,}75\,;\,-1{,}28\,;\,-0{,}84\,;\,
0{,}84\,;\,1{,}28\,;\,1{,}75\}$. \\
\bottomrule
\end{tabularx}
\end{center}

Deux décisions de conception méritent d'être signalées.

\paragraph{Une référence qui se dégrade explicitement.} Toutes les stations ne disposent pas de
trente années d'historique. Plutôt que d'écarter les stations courtes ou de leur appliquer
silencieusement une référence plus faible, la sélection de fenêtre se dégrade par paliers
déclarés : la normale 1991--2020 si au moins quinze années sont disponibles, sinon la meilleure
fenêtre trentenaire alignée sur les décennies, sinon l'historique complet. Le palier retenu est
inscrit dans une colonne accompagnant chaque valeur, de sorte qu'un utilisateur sait toujours
sur quelle base une station est classée.

\paragraph{Une grille n'est jamais fabriquée.} Si aucun mois n'est exploitable, la grille reste
nulle et l'indice prend la valeur \emph{inconnu}. Produire une classe plutôt qu'un aveu
d'ignorance reviendrait à inventer une information que l'utilisateur lirait comme une mesure.

La table de référence est recalculée \emph{hebdomadairement}, et non chaque nuit ni une fois par
décennie : la fenêtre est fixe, mais les grilles par station ne le sont pas — une station
accumule de l'historique, de nouvelles stations apparaissent, et une station classée en
référence provisoire peut franchir le seuil des quinze années.

\section{Les indices climatiques sur grille}

Trois indices sont calculés pour chaque maille de \ang{0.1}, sur quatre fenêtres d'agrégation de
1, 3, 6 et 12 mois correspondant à des inerties différentes du système hydrologique. Tous
emploient la période de référence 1991--2020 et l'échelle à sept classes de
McKee~\cite{mckee1993spi}, adoptée par l'Organisation météorologique mondiale.

\begin{center}
\begin{tabularx}{\textwidth}{@{}llX@{}}
\toprule
\textbf{Indice} & \textbf{Grandeur} & \textbf{Ajustement} \\
\midrule
SPI  & Précipitation & Loi gamma, mélangée à la probabilité de cumul nul \\
STI  & Température & Loi normale (score centré réduit) \\
SPEI & Précipitation moins évapotranspiration & Loi logistique généralisée \\
\bottomrule
\end{tabularx}
\end{center}

\section{Premier résultat : l'évapotranspiration de référence a dû être recalculée}
\label{sec:hargreaves}

Le calcul du SPEI suppose un bilan hydrique, donc une évapotranspiration. ERA5-Land fournit
directement une variable d'évaporation potentielle : l'employer semblait naturel, et c'est ce
qui a été fait initialement.

Le résultat était aberrant. Comparé sur \num{30888} couples maille-mois entre 2015 et 2025, sur
les mêmes mailles :

\begin{center}\small
\begin{tabularx}{\textwidth}{@{}Xrrr@{}}
\toprule
 & \textbf{Hargreaves (FAO-56)} & \textbf{Évaporation potentielle ERA5-Land} & \textbf{Rapport} \\
\midrule
Évapotranspiration annuelle & \SI{818}{\milli\metre} & \SI{1756}{\milli\metre} & $\times\,\num{2,15}$ \\
Bilan précipitation $-$ ETP & \SI{+146}{\milli\metre\per\an} & \SI{-793}{\milli\metre\per\an} & --- \\
\bottomrule
\end{tabularx}
\end{center}

La seconde colonne place \textbf{l'intégralité du territoire français en déficit hydrique
permanent}, y compris les années humides — ce qui est physiquement faux. La valeur de
\SI{818}{\milli\metre\per\an} se situe dans l'intervalle des références publiées pour la
France (\SIrange{700}{900}{\milli\metre\per\an}) ; celle de \SI{1756}{\milli\metre\per\an}
n'y est pas.

L'explication tient à la définition des grandeurs, et c'est le point à retenir : les deux
portent des noms voisins et ne décrivent pas la même chose. L'évaporation potentielle d'ERA5
n'est pas une évapotranspiration de référence au sens de la FAO — c'est l'évaporation d'une
surface non contrainte en eau, calculée avec la résistance aérodynamique du modèle
atmosphérique, dont la littérature documente la surestimation substantielle. La formule de
Hargreaves~\cite{hargreaves1985} est précisément le repli recommandé par la
FAO-56~\cite{allen1998fao56} lorsque le rayonnement, le vent et l'humidité ne sont pas
disponibles ; elle ne requiert que les températures extrêmes journalières et le rayonnement
extraterrestre, dérivable analytiquement de la latitude et du jour de l'année. Elle est
implémentée directement en SQL, avec le calcul complet du rayonnement extraterrestre.

\begin{resultat}[title={Décision}]
L'évapotranspiration de référence est recalculée par la formule de Hargreaves à partir des
températures journalières vraies. La variable ERA5 d'origine reste exposée pour la traçabilité,
avec une consigne explicite de non-consommation.
\end{resultat}

\section{Deuxième résultat : le choix de la loi statistique du SPEI}
\label{sec:glo}

La formulation de référence du SPEI~\cite{vicenteserrano2010spei} ajuste une loi log-logistique
à trois paramètres. Appliquée à la grille française, elle ne converge que sur
\SI{74,6}{\percent} des triplets maille~$\times$~mois~$\times$~fenêtre. Un quart du domaine
restait sans valeur, ce qui est inacceptable pour une couche cartographique nationale.

Plutôt que de constater le taux d'échec, les rejets ont été instrumentés pour en identifier la
cause. Le résultat est net : \textbf{la totalité des rejets} provient de la garde sur le
paramètre de forme, et \textbf{aucun} d'une insuffisance de données.

L'explication tient en trois phrases. L'ajustement passe par la méthode des
L-moments~\cite{hosking1990lmoments}, qui résume la distribution par un coefficient d'asymétrie
noté $\tau_3$. La loi log-logistique exige une asymétrie positive. Or environ
\SI{27}{\percent} des mailles présentent une asymétrie négative — une distribution du bilan
hydrique dont la queue s'étend vers les valeurs basses — et aucune taille d'échantillon ne
corrige cela : la loi ne peut structurellement pas représenter cette forme.

La logistique généralisée lève la limite : son paramètre de forme vaut $k = -\tau_3$ et accepte
donc les deux signes. C'est également la loi retenue par l'implémentation de référence du
SPEI~\cite{begueria_spei_package}. La couverture atteint \SI{100}{\percent}.

\begin{resultat}[title={Un changement additif, non destructif}]
La log-logistique est exactement la logistique généralisée reparamétrée. Le remplacement laisse
donc les valeurs inchangées partout où la loi d'origine convergeait — vérifié à un écart maximal
de \num{0,000} sur \num{35614} mailles — et ne fait qu'ajouter les valeurs manquantes.
\end{resultat}

\section{Troisième résultat : la validation des indices produits}

Un indice standardisé doit, par construction, suivre approximativement une loi normale centrée
réduite \emph{sur sa propre période de référence}. C'est une propriété vérifiable, et elle a été
vérifiée.

\paragraph{Protocole.} Le critère d'acceptation a été fixé \textbf{avant} la mesure : moyenne
voisine de zéro à $\pm\,\num{0,05}$, écart-type voisin de un à $\pm\,\num{0,05}$, et taux de
saturation aux bornes inférieur à \SI{1}{\percent}. Fixer le critère avant de mesurer évite
l'ajustement rétrospectif du seuil au résultat obtenu.

\paragraph{Mesure.} Effectuée sur \num{41960400} valeurs couvrant \num{11496} mailles.

\begin{center}
\begin{tabularx}{\textwidth}{@{}lXrrrr@{}}
\toprule
\textbf{Indice} & \textbf{Fenêtres} & \textbf{Moyenne} & \textbf{Écart-type} & \textbf{Médiane}
& \textbf{Saturation} \\
\midrule
SPI  & 1/3/6/12 & $-0{,}008$ à $+0{,}002$ & \numrange{0,985}{1,031} & \numrange{0,00}{0,04} & \SIrange{0,05}{0,51}{\percent} \\
STI  & 1/3/6/12 & $0{,}000$               & \num{0,983}            & $-0{,}09$ à $-0{,}02$ & \SIrange{0,00}{0,04}{\percent} \\
SPEI & 1/3/6/12 & $+0{,}004$ à $+0{,}006$ & \numrange{0,999}{1,013} & \numrange{0,01}{0,03} & \SIrange{0,012}{0,035}{\percent} \\
\bottomrule
\end{tabularx}
\end{center}

\begin{resultat}[title={Conclusion}]
Les trois indices satisfont le critère d'acceptation sur les quatre fenêtres. La couverture du
SPEI est uniforme à \SI{100}{\percent}, ce qui signifie que la calibration n'est pas seulement
préservée mais mesurée sur l'ensemble du domaine.
\end{resultat}

\section{Deux lectures à ne pas faire}

La validation ci-dessus autorise l'usage des indices ; elle n'autorise pas toutes les
interprétations. Deux mises en garde ont été consignées, parce que l'erreur correspondante est
naturelle.

\paragraph{La classe « normale » ne couvre pas la proportion théorique.} Elle représente
\SI{54,8}{\percent} des valeurs au lieu des \SI{59,9}{\percent} attendus d'une gaussienne, alors
même que l'écart-type vaut \num{1,008}. La distribution conserve des épaules légèrement plus
lourdes. Cet écart est une propriété de la transformation au voisinage des bornes de classe. Il
n'a pas de conséquence opérationnelle — les seuils employés sont ceux, communs, des autres
indices — et il ne doit surtout \textbf{pas} être lu comme une anomalie de sécheresse.

\paragraph{Une décroissance décennale est un signal climatique, non une dérive de méthode.} La
moyenne du SPEI sur fenêtre d'un mois décroît de manière monotone d'une décennie à l'autre :
$+\num{0,038}$ dans les années 1990, $+\num{0,007}$ dans les années 2000, $-\num{0,016}$ dans les
années 2010, $-\num{0,117}$ depuis 2020. Cette évolution avait été anticipée avant d'être
mesurée, ce qui la distingue d'un artefact d'ajustement : elle traduit un assèchement, non un
biais.

\section{Un contrat entre les deux dépôts}

La classification en sept classes existe en deux endroits : dans l'entrepôt, qui produit les
indicateurs, et dans la plateforme, qui les affiche et les compare aux prévisions. Cette
duplication est délibérée — elle évite un couplage fort entre deux systèmes déployés séparément
— mais elle crée un risque de divergence silencieuse.

Le risque est traité par un \textbf{contrat testé} : chaque dépôt possède une table de référence
associant des valeurs d'entrée aux classes attendues, et ces deux tables sont identiques par
construction. Toute modification de l'une sans l'autre fait échouer une suite de tests. Le
mécanisme est simple, et c'est sa vertu : il rend la contrainte exécutable plutôt que
documentaire.

\begin{acquisouvert}
\textbf{Acquis} — Une famille d'indicateurs normalisés, produite à l'échelle nationale, validée
selon un protocole à critère pré-établi, et dont deux décisions méthodologiques structurantes
sont documentées avec leurs mesures. Ces indicateurs sont directement utilisables par les
hydrogéologues et constituent des variables candidates pour la modélisation.

\textbf{Ouvert} — Une incohérence connue subsiste entre la chaîne « grille » et la chaîne
« station » sur l'évapotranspiration : la seconde expose encore la variable d'origine, car elle
constitue l'entrée de forçage des modèles métier déjà calibrés et sa modification les
invaliderait. Cet arbitrage est explicite et documenté ; sa résolution suppose une
recalibration d'ensemble.
\end{acquisouvert}
